%% kapa: exam=1 round=31 year=2020 period=2 form=A subject=accounting subject_ko=회계학 seq=18 number=58
%% source: Q-Net 감정평가사 기출문제 (cst00309) artlSeq=5207219 file=extracted/1008/5207219/A형/A형/2020년도 제31회 감정평가사 1차 시험_2교시_A형.pdf
%% conversion: kordoc 4.24.51 -> tools/gen1.py (auto); review=proofread-sonnet; flags=none
\begin{question}{58}
(주)감평은 20x1년 초 기계장치(취득원가 ￦1,000,000, 내용연수 5년, 잔존가치 ￦0, 정액법 상각)를 취득하여 원가모형을 적용하고 있다. 20x2년 초 (주)감평은 동 기계장치에 대해 자산인식기준을 충족하는 후속원가 ￦325,000을 지출하였다. 이로 인해 내용연수가 2년 연장(20x2년 초 기준 잔존내용연수 6년)되고 잔존가치는 ￦75,000 증가할 것으로 추정하였으며, 감가상각방법은 이중체감법(상각률은 정액법 상각률의 2배)으로 변경하였다. (주)감평은 동 기계장치를 20x3년 초 현금을 받고 처분하였으며, 처분이익은 ￦10,000이다. 기계장치 처분 시 수취한 현금은?
\begin{choices}
  \item ￦610,000
  \item ￦628,750
  \item ￦676,667
  \item ￦760,000
  \item ￦785,000
\end{choices}
\end{question}
\answer{4}
